\section{Doubling and separating interfaces}\label{sec:doubling-bubbles}

Fix a component $P$ of the cut manifold in
Proposition~\ref{prop:separating-cut}.  We will construct small separating
surfaces between a base point and a distant distance slab.  The boundary
of $P$ prevents a direct use of the closed-surface estimate, so we first
smooth the doubled weighted geometry.  We then minimize weighted
perimeter with a pressure that forces the resulting surfaces to lie in
a prescribed annular region.

\subsection{Smoothing the double}

Let $Z$ be the double of $P$ along $\partial P$, and let $\sigma$ exchange
its two copies.  If $\partial P$ is empty, take two disjoint copies.

\begin{lemma}[Weighted doubling]\label{lem:double-smoothing}
There are smooth $\sigma$-invariant data $(h,T)$ on $Z$ such that
\begin{equation}\label{eq:double-comparison}
 D(h,T)\ge\frac12,
 \qquad \frac14g\le h\le4g
 \quad\text{on each copy of }P.
\end{equation}
The data can be left unchanged outside disjoint, locally finite collars
of the boundary components.
\end{lemma}

\begin{proof}
Choose disjoint normal collars of the compact boundary components.  Their
widths may vary with the component.  In one such collar, use inward
distance $s\ge0$ and write
\begin{equation}\label{eq:weighted-collar-data}
 g=ds^2+k(s),\qquad
 m(s)=\log\frac{d\vol_{k(s)}}{d\vol_{k(0)}}-t(s).
\end{equation}
The ratio in this formula is a ratio of densities on the boundary
surface, so no orientation is needed.  With the normal $\partial_s$,
the slices have mean curvature
$H=\frac12\operatorname{tr}_{k(s)}k'(s)$, and
$m'=H-t'$.  Weighted minimality of the boundary therefore gives
\begin{equation}\label{eq:even-weighted-density}
 m'(0)=0.
\end{equation}
The useful curvature formula in these variables is
\begin{equation}\label{eq:collar-curvature}
 D(g,t)=D(k(s),t(s))
       -\left|\frac12k'(s)\right|_{k(s)}^2-(m')^2-2m''.
\end{equation}
Here $D(k(s),t(s))$ uses the two-dimensional slice metric.  To verify the
formula, use
\[
 \Scal_g=\Scal_{k(s)}-\left|\frac12k'\right|^2-H^2-2H',
 \qquad
 \Delta_g t=\Delta_{k(s)}t+t''+Ht',
\]
and combine the normal terms as
$-(H-t')^2-2(H'-t'')$.

Reflect $k$ and $m$ evenly across $s=0$.  The reflected $k$ is continuous
and locally Lipschitz, whereas the reflected $m$ is $C^2$ by
\eqref{eq:even-weighted-density}.  Convolve these data in $s$ with an even,
nonnegative smooth mollifier of radius $\delta$.  An even cutoff, equal
to one near zero and supported in a smaller collar, blends the
convolution back to the original data.  Denote the resulting smooth
even data by $k_\delta,m_\delta$, and set
\[
 h_\delta=ds^2+k_\delta(s),\qquad
 T_\delta=\log\frac{d\vol_{k_\delta(s)}}{d\vol_{k(0)}}-m_\delta(s).
\]
For small $\delta$, the tensor $k_\delta$ is positive definite.

We check the curvature inequality also at points approaching the seam.
On compact subcollars, $k_\delta,m_\delta$ and their tangential derivatives
through order two converge uniformly to the reflected data.  The first
two normal derivatives of $m_\delta$ converge uniformly as well.  Thus
every term in \eqref{eq:collar-curvature}, except the squared norm of
$k_\delta'/2$, converges to its one-sided value at the seam.  For a
sequence of points approaching $(x,0)$ as $\delta\to0$, the convolved
$k'$ averages tensors approaching $k'(x,0)$ and $-k'(x,0)$.  Convexity of
the squared norm, and convergence of the metrics used in that norm,
give
\[
 \limsup_{\delta\to0}
 \left|\frac12k_\delta'\right|_{k_\delta}^2
 \le \left|\frac12k'(x,0)\right|_{k(0)}^2.
\]
Consequently the lower limit of $D(h_\delta,T_\delta)$ there is at least
one.  Away from the seam, the convergence is smooth.  Compactness of
the collar region in which changes occur now gives
$D(h_\delta,T_\delta)\ge1/2$ for all sufficiently small $\delta$.
Uniform convergence also gives the metric comparison in
\eqref{eq:double-comparison}.  Choose the smoothing radius separately
on each collar.  The collars are locally finite, so these choices give
global smooth data, equal to $(g,t)$ elsewhere.
\end{proof}

Write $d$ for the intrinsic length metric of $h|_P$.  It is proper by
\eqref{eq:double-comparison} and properness of the original length metric
on $P$.  Fix $o\in\operatorname{int}P$ and put
\begin{equation}\label{eq:radial-distance}
 r(x)=d(o,x)\qquad(x\in P).
\end{equation}
Extend $r$ evenly to $Z$.  The folding map $Z\to P$ preserves the length
of a path on either copy and does not increase length across the seam.
This follows directly from $h=ds^2+k(s)$, with $k$ even, by replacing
$s$ with $|s|$ along an absolutely continuous path.  In particular, the
even extension of $r$ is proper and $1$-Lipschitz on each component of
$Z$.  We will also use the consequence
\begin{equation}\label{eq:folding-distance}
 d(x,y)\le\operatorname{length}_h(\gamma)
 \quad(x,y\in P)
\end{equation}
for every path $\gamma$ in $Z$ from $x$ to $y$.

\subsection{Separating interfaces in a distance annulus}

We use a weighted perimeter-minus-volume minimization, as in
\cite[Proposition~12]{ChodoshLi2024}, to obtain the following output for
the slab argument.  An open set here is understood relative to $P$ when
it meets $\partial P$.

\begin{proposition}[Separating interfaces]\label{prop:separating-interfaces}
Let $\ell\ge100$, and suppose that $\{r>\ell\}$ is nonempty.  There are finitely
many disjoint smooth compact connected two-sided neat hypersurfaces
$F_1,\ldots,F_m$ in $P$ such that
\begin{equation}\label{eq:interface-location}
 F_i\subset\{\ell-91<r<\ell-18\},\qquad \diam_d F_i\le4\pi.
\end{equation}
Each $F_i$ separates $P$.  Moreover,
$P\setminus\bigcup_iF_i$ is the disjoint union of two open sets
$U_+,U_-$, with $o\in U_+$ and $\{r>\ell\}\subset U_-$.
\end{proposition}

\begin{proof}
First choose a smooth even function $\rho$ on $Z$ such that
\begin{equation}\label{eq:smooth-radial-distance}
 |\rho-r|<1,\qquad |d\rho|_h\le2.
\end{equation}
For completeness, smooth a Lipschitz function in relatively compact
coordinate charts in which the metric is sufficiently close to a
constant metric.  Local convolution then gives gradient norm at most
$3/2$, with any prescribed positive value error on the compact support
of a partition function.  Let $f_j$ be these local smoothings and patch
them using a countable, locally finite smooth partition of unity
$\{\phi_j\}$.  Since
$\sum_jd\phi_j=0$, the error in the gradient caused by patching is
$\sum_j(f_j-r)d\phi_j$.  Taking the value errors smaller than
$2^{-j-2}/(1+\|d\phi_j\|_\infty)$ makes its norm less than $1/2$ and
also makes the total value error less than one.  Finally average the
result with its reflection.  Since reflection is an isometry, both
bounds survive.  The function $\rho$ has compact sublevel sets because
$r$ does.

Choose regular values
\begin{equation}\label{eq:pressure-levels}
 a\in(\ell-90,\ell-89),\qquad b\in(\ell-20,\ell-19).
\end{equation}
On the band $a<\rho<b$ define
\begin{equation}\label{eq:cotangent-pressure}
 \mu=\cot\left(\frac{\pi(\rho-a)}{b-a}\right).
\end{equation}
The width $b-a>69$ and \eqref{eq:smooth-radial-distance} imply
\begin{equation}\label{eq:pressure-gradient}
 2|d\mu|_h
 \le\frac{4\pi}{b-a}(1+\mu^2)
 \le\frac14(1+\mu^2).
\end{equation}
The pressure tends to $+\infty$ at the inner end of the band and to
$-\infty$ at the outer end.  These signs will keep a minimizing
interface away from the constraints used to obtain it.

\subsubsection*{Existence and reflection symmetry}
For sufficiently small $\epsilon>0$, minimize among sets $E$ of finite
perimeter satisfying, up to sets of measure zero,
\begin{equation}\label{eq:phase-constraints}
 \{\rho<a+\epsilon\}\subset E,
 \qquad E\subset\{\rho\le b-\epsilon\}.
\end{equation}
Write
\[
 \mathcal P_T(E)=\int_Z e^{-T}\,d|D\chi_E|_h,
 \qquad
 \mathcal J(E)=\mathcal P_T(E)
       -\int_E\mu_*e^{-T}\,d\vol_h,
\]
where $\mu_*$ is smooth, compactly supported, reflection invariant, and
agrees with $\mu$ on a neighborhood of
$\{a+\epsilon\le\rho\le b-\epsilon\}$.  Such an extension exists
because this closed band is compactly contained in $\{a<\rho<b\}$.
The perimeter is taken on all of $Z$, not relative to the band.

Regular sublevel sets provide competitors.  All admissible sets lie
in a fixed compact set, where the smooth weight is bounded above and
below by positive constants.  A minimizing sequence therefore has
uniformly bounded ordinary perimeter; compactness for functions of
bounded variation gives an $L^1$-convergent subsequence.  The constraints
are closed under this convergence, the volume term converges, and
weighted perimeter is lower semicontinuous; see, for example,
\cite{Maggi2012}.  Thus a minimizer exists.

The lattice inequality for perimeter gives
\begin{equation}\label{eq:perimeter-lattice}
 \mathcal J(E\cup E')+\mathcal J(E\cap E')
 \le\mathcal J(E)+\mathcal J(E').
\end{equation}
Indeed the volume terms satisfy equality, while perimeter satisfies
the indicated inequality.  The reflected set $\sigma E$ is another
minimizer, and both its union and intersection with $E$ are admissible.
Minimality and \eqref{eq:perimeter-lattice} show that both are minimizers.
Replacing $E$ by $E\cup\sigma E$, we may suppose that $E$ is invariant
under reflection, up to a null set.

\subsubsection*{Removing the constraints near the interface}
We show that $E$ fills a larger inner sublevel and avoids a larger
outer superlevel.  This is the step that permits unrestricted local
variations of its boundary.

Since $a$ is a regular value and its level set is compact, there is a
fixed compact level neighborhood on which $d\rho\ne0$.  Choose a smooth
compactly supported field $X$, with $|X|_h\le1$, that equals
$\nabla\rho/|d\rho|_h$ on a smaller such neighborhood.  This one field
can be used for all sufficiently small $\epsilon$.  Set
$A=\{\rho<a+2\epsilon\}$.  The field $X$ equals the outward unit normal
along $\partial A$, so the divergence theorem and the variational
definition of perimeter give
\begin{align}
 \mathcal P_T(E\cup A)-\mathcal P_T(E)
 &\le\mathcal P_T(A)-\mathcal P_T(E\cap A)\notag\\
 &\le\int_{A\setminus E}\Div_h(e^{-T}X)\,d\vol_h.
 \label{eq:inner-barrier-comparison}
\end{align}
By \eqref{eq:phase-constraints}, the difference set is contained, up
to a null set, in $\{a+\epsilon\le\rho<a+2\epsilon\}$.  On this shell,
$\mu_*e^{-T}>\Div_h(e^{-T}X)$ when $\epsilon$ is sufficiently small:
the divergence is bounded on a fixed compact set, while the pressure
tends uniformly to $+\infty$.  Subtracting the volume terms in
\eqref{eq:inner-barrier-comparison} shows that $E\cup A$ would strictly
improve $\mathcal J$ unless $A\setminus E$ has measure zero.  Hence
\begin{equation}\label{eq:inner-phase-clearance}
 \{\rho<a+2\epsilon\}\subset E
 \quad\text{up to a null set}.
\end{equation}

For the outer end apply the same comparison to the complement of $E$,
whose pressure is $-\mu_*$.  Use
$B=\{\rho>b-2\epsilon\}$ and a field equal to
$-\nabla\rho/|d\rho|_h$ near its boundary.  Both the field and the
change of phase have compact support near that boundary; thus the
divergence calculation also applies if $B$ itself is unbounded.
The pressure $-\mu_*$ tends uniformly to $+\infty$ on
$\{b-2\epsilon<\rho\le b-\epsilon\}$.  It follows that
\begin{equation}\label{eq:outer-phase-clearance}
 E\cap\{\rho>b-2\epsilon\}=\varnothing
 \quad\text{up to a null set}.
\end{equation}

Let $S=\supp|D\chi_E|_h$.  Equations
\eqref{eq:inner-phase-clearance} and
\eqref{eq:outer-phase-clearance} show that $S$ is a compact subset of
the open band $a+\epsilon<\rho<b-\epsilon$.  Every compactly supported
change of $E$ in this band is admissible.  The minimizer therefore has
the full local minimizing property for weighted perimeter with the
smooth pressure $\mu$ near every point of $S$.

\subsubsection*{Regularity, diameter, and restriction to the half}
We use the following precise local regularity input: in a smooth
three-manifold, a set locally minimizing perimeter with a smooth
positive weight and a smooth pressure volume term has smooth embedded
perimeter support, and its outward measure-theoretic unit normal
extends smoothly to that support.  For the weighted version, see
\cite[Theorem~4.3 and Remarks~4.4]{Kemper2020}: the singular-set bound is
ambient dimension minus eight, and smooth data give smooth regular
parts by elliptic regularity.  The weight can equivalently be absorbed
in the smooth conformal metric $\bar h=e^{-T}h$: in dimension three its
surface area is the weighted $h$-area, and its pressure is
$\mu e^{T/2}$.  Only this local regularity theorem is used here, with
the full-neighborhood minimizing property just established.
In particular $S$ is a compact smooth two-sided hypersurface with
finitely many connected components.  On its complement the phase is
locally constant: a characteristic function with zero distributional
derivative on a connected open set is constant almost everywhere.

Each connected component $\Sigma$ of $S$ is stationary and stable for
all normal variations of the weighted perimeter-minus-volume
functional.  Its outward normal satisfies
\[
 H-\nu T=\mu.
\]
From \eqref{eq:double-comparison} and \eqref{eq:pressure-gradient},
\[
 D(h,T)+\mu^2-2|d\mu|_h
 \ge\frac14+\frac34\mu^2\ge\frac14.
\]
Proposition~\ref{prop:stable-surface}, with $c=1/4$, now gives
\begin{equation}\label{eq:closed-interface-diameter}
 \diam_{h|_\Sigma}\Sigma\le4\pi.
\end{equation}

Reflection preserves the full and empty phases and hence $S$ with its
outward normal.  At a point of $S$ on the seam, that normal is fixed by
$d\sigma$, so it is tangent
to the seam.  Thus $S$ meets the seam transversally.  Its restriction
to $P$ is consequently a finite union of disjoint compact two-sided
neat hypersurfaces.  Let $F_1,\ldots,F_m$ be their connected components.
Two points of one $F_i$ lie in one connected component $\Sigma$ of
$S$.  A path in $\Sigma$ of length at most $4\pi$, followed by folding,
and \eqref{eq:folding-distance} give $\diam_d F_i\le4\pi$.
Proposition~\ref{prop:separating-cut} says that every $F_i$ separates
$P$.

The location of $S$ in $a<\rho<b$, together with
\eqref{eq:smooth-radial-distance} and \eqref{eq:pressure-levels}, gives
\eqref{eq:interface-location}.  Restrict the full and empty phases of
$Z\setminus S$ to $P$ to obtain $U_+,U_-$.  They are disjoint relatively
open sets covering $P\setminus\bigcup_iF_i$.  Since $\rho(o)<1<a$,
the base point is in the full phase.  If $r(x)>\ell$, then
$\rho(x)>\ell-1>b$, so $x$ is in the empty phase.  This proves the final
assertion.
\end{proof}
